\chapter{Réalisation et résultats}
\label{chap:realisation}

\section{Introduction}
Ce chapitre présente l'environnement de développement et les interfaces réalisées, puis l'évaluation chiffrée du prototype. Il se termine par une analyse honnête de ses limites. Toutes les captures d'écran ont été prises sur l'application en fonctionnement.

\section{Environnement technique}

\begin{table}[H]
\centering
\caption{Technologies utilisées}
\label{tab:techno}
\begin{tabularx}{\textwidth}{l X}
\toprule
\textbf{Domaine} & \textbf{Outils} \\
\midrule
Acquisition & KEPServerEX V6 (OPC UA)~\cite{kepware}, Node-RED et \texttt{node-red-contrib-opcua}~\cite{nodered} \\
Transport & Eclipse Mosquitto 2.0 (MQTT)~\cite{mqtt5}, paho-mqtt \\
Back-end & Python, FastAPI~\cite{fastapi}, pandas, PyJWT, bcrypt \\
Stockage & SQLite~\cite{sqlite}, Neo4j 5~\cite{neo4j}, fichiers CSV \\
Apprentissage & scikit-learn~\cite{sklearn}, XGBoost~\cite{chen2016}, joblib \\
IA générative & API Google Gemini, Graph RAG~\cite{edge2024} \\
Front-end & React 19~\cite{react}, Vite, React Three Fiber~\cite{r3f}, vis-network, Recharts \\
Déploiement & Docker Compose~\cite{merkel2014} : 5 services (broker, Neo4j, backend, publieur, frontend) \\
\bottomrule
\end{tabularx}
\end{table}

\section{Interfaces réalisées}

\subsection{Authentification}
L'accès passe par une page de connexion (figure~\ref{fig:login}). Le backend vérifie le mot de passe haché avec bcrypt, puis délivre un jeton \gls{jwt} valable 8 heures. Chaque page de l'application est protégée : sans jeton valide, l'utilisateur est renvoyé vers la connexion.

\begin{figure}[H]
\centering
\includegraphics[width=0.8\textwidth]{app_login}
\caption{Page de connexion, avec les trois profils de démonstration}
\label{fig:login}
\end{figure}

\subsection{Synoptique et graphe de connaissances}
La vue principale juxtapose le schéma du circuit et le graphe de connaissances (figure~\ref{fig:home}). Une barre de pilotage permet de lancer la relecture et de choisir sa vitesse. Elle donne aussi accès aux panneaux de chaque source de données et au simulateur What-If.

\begin{figure}[H]
\centering
\includegraphics[width=\textwidth]{app_home}
\caption{Vue double : synoptique du circuit (à gauche) et graphe de connaissances (à droite)}
\label{fig:home}
\end{figure}

Pendant la relecture, les lignes s'animent selon les débits reçus, et chaque équipement s'illumine selon son état (figure~\ref{fig:flowsheet}).

\begin{figure}[H]
\centering
\includegraphics[width=\textwidth]{app_flowsheet_full}
\caption{Synoptique en plein écran pendant la relecture}
\label{fig:flowsheet}
\end{figure}

\subsection{Fiche d'un équipement}
Un clic sur un équipement ouvre sa fiche (figure~\ref{fig:drawer}). Pour le broyeur BM\_001, elle affiche :
\begin{itemize}
  \item la télémétrie en direct : puissance, courant, vitesse, températures des paliers, vibration, niveau sonore ;
  \item les indicateurs dérivés calculés par le backend : réduction de taille, débits entrant et sortant, écart de débit.
\end{itemize}
Le graphe se centre alors automatiquement sur cet équipement et met en évidence ses liens : capteurs (\texttt{MONITORS}), lignes (\texttt{FEEDS\_INTO}, \texttt{DISCHARGES\_TO}) et risques (\texttt{RISKS}).

\begin{figure}[H]
\centering
\includegraphics[width=\textwidth]{app_drawer_bm}
\caption{Fiche du broyeur BM\_001 : télémétrie en direct, indicateurs dérivés et voisinage dans le graphe}
\label{fig:drawer}
\end{figure}

Le graphe de connaissances compte \textbf{69 nœuds et 129 relations}, répartis ainsi :
\begin{itemize}
  \item 6 équipements, 7 lignes et 15 capteurs ;
  \item 6 modes de défaillance ;
  \item 29 bons de travail et 6 techniciens.
\end{itemize}
Des filtres isolent chaque famille de nœuds, et un mode de parcours trace le chemin d'une cause à ses effets.

\subsection{Simulateur What-If}
Le panneau What-If (figure~\ref{fig:whatif}) propose un curseur pour chacune des 8 entrées réglables. À chaque modification, il affiche pour chaque grandeur aval :
\begin{itemize}
  \item la valeur actuelle et la valeur prédite ;
  \item l'écart entre les deux ;
  \item le $R^2$ du modèle sur cette grandeur. Un badge orange signale une prédiction moins fiable, par exemple sur les pressions des cyclones.
\end{itemize}

\begin{figure}[H]
\centering
\includegraphics[width=\textwidth]{app_whatif}
\caption{Simulateur What-If : entrées réglables et prédictions comparées à l'état actuel}
\label{fig:whatif}
\end{figure}

\subsection{Vue 3D}
La vue 3D (figure~\ref{fig:3d}) reproduit l'implantation du circuit : bac, pompe, batterie de trois hydrocyclones, broyeur et tuyauteries. Au-dessus de chaque équipement, une étiquette affiche sa grandeur clé en direct : niveau du bac, pression de refoulement de la pompe, débit d'alimentation des cyclones, puissance du broyeur.

\begin{figure}[H]
\centering
\includegraphics[width=\textwidth]{app_view3d}
\caption{Vue 3D du circuit réalisée avec React Three Fiber}
\label{fig:3d}
\end{figure}

\subsection{Maintenance et fiabilité}
La page de maintenance (figure~\ref{fig:maint}) résume l'état du parc : nombre d'équipements, interventions en retard, équipements critiques et MTBF moyen. Elle donne aussi accès à chaque équipement, avec le nombre de jours restant avant sa prochaine intervention. La fiche détaillée (figure~\ref{fig:maintpump}) présente :
\begin{itemize}
  \item les données constructeur de l'équipement ;
  \item la dernière intervention et la prochaine échéance ;
  \item l'état de l'équipement, avec son MTBF et son MTTR ;
  \item l'état de ses composants critiques.
\end{itemize}
Un lien ouvre l'historique complet des interventions (31 bons de travail au total).

\begin{figure}[H]
\centering
\includegraphics[width=\textwidth]{app_maintenance}
\caption{Tableau de bord de maintenance et de fiabilité}
\label{fig:maint}
\end{figure}

\begin{figure}[H]
\centering
\includegraphics[width=\textwidth]{app_maintenance_pump}
\caption{Fiche de maintenance de la pompe SP\_001 (criticité haute, revêtement d'usure à remplacer)}
\label{fig:maintpump}
\end{figure}

\subsection{Copilote IA}
Le copilote (figure~\ref{fig:copilot}) est une fenêtre flottante disponible sur toutes les pages. Il propose des questions types, et chaque réponse indique le moteur utilisé et le nombre d'entités du graphe mobilisées. Dans l'exemple ci-dessous, pris hors connexion à l'API Gemini, le moteur local a construit sa réponse à partir de 15 entités. On y retrouve la vitesse du broyeur (74,05\,\% de la vitesse critique), sa puissance (1\,818~kW) et la réduction de taille réalisée.

\begin{figure}[H]
\centering
\includegraphics[width=\textwidth]{app_copilot}
\caption{Copilote Graph RAG : réponse sur l'état du broyeur BM\_001}
\label{fig:copilot}
\end{figure}

\section{Évaluation}

\subsection{Capteur virtuel du P80 par XGBoost}
\label{sec:xgboost}

\paragraph{Objectif.} Estimer en continu le P80 de la surverse des hydrocyclones, c'est-à-dire la finesse du produit envoyé en aval, sans attendre une analyse granulométrique.

\paragraph{Données.} Le modèle est entraîné sur les données du jumeau (dossier \texttt{data/}) : 2\,880 mesures procédé acquises toutes les 15 minutes du 1\textsuperscript{er} au 30 juillet 2026, et 43\,200 mesures de santé acquises chaque minute. Chaque mesure procédé est associée à la \textbf{moyenne des signaux de santé sur les 15 minutes qui la précèdent}, ce qui reprend le principe d'alignement temporel décrit à la section~\ref{sec:asymetrie}. Sur cette période, le P80 de la surverse varie entre 135,7 et 149,9~µm (moyenne 142,7~µm, écart-type 3,67~µm).

\paragraph{Choix des variables et prévention des fuites.} Un capteur virtuel ne doit utiliser que des grandeurs réellement disponibles au moment de l'estimation. Toutes les granulométries et teneurs mesurées en aval, ainsi que le rapport de réduction (calculé à partir du P80), ont donc été \textbf{exclues}, car elles contiennent la réponse. Deux scénarios ont été évalués :
\begin{itemize}
  \item \textbf{Scénario A -- capteurs en ligne} (32 variables) : débits et fractions solides des différentes lignes (débitmètres et densimètres), ainsi que les signaux de santé des équipements (puissance et vitesse du broyeur, pressions des cyclones, vitesse de la pompe, niveau du bac, etc.) ;
  \item \textbf{Scénario B -- A + caractérisation de l'alimentation} (34 variables) : on ajoute le P80 et la teneur BPL de l'alimentation fraîche, connus en amont du circuit.
\end{itemize}

\paragraph{Protocole.} Pour respecter la nature temporelle des données, le découpage est \textbf{chronologique} : 70\,\% pour l'entraînement (2\,015 points), 15\,\% pour la validation (433 points) et 15\,\% pour le test (432 points, du 26 au 30 juillet). Un découpage aléatoire aurait laissé le modèle «~voir le futur~» et surestimé ses performances. Les hyperparamètres de XGBoost~\cite{chen2016} (profondeur, taux d'apprentissage, sous-échantillonnage) sont choisis par recherche de grille, avec arrêt anticipé sur le jeu de validation. Le modèle final est ensuite réentraîné sur l'entraînement et la validation, puis évalué une seule fois sur le test. Trois modèles servent de comparaison :
\begin{itemize}
  \item une prédiction constante égale à la moyenne, qui sert de plancher ;
  \item une régression linéaire régularisée (Ridge) ;
  \item un Random Forest avec les mêmes hyperparamètres que le simulateur What-If du backend.
\end{itemize}
Les performances sont mesurées par le coefficient de détermination $R^2$, la \gls{rmse}, la \gls{mae} et l'erreur relative moyenne (MAPE) :
\[
\mathrm{RMSE}=\sqrt{\frac{1}{n}\sum_{i=1}^{n}(y_i-\hat{y}_i)^2}, \qquad
\mathrm{MAE}=\frac{1}{n}\sum_{i=1}^{n}\lvert y_i-\hat{y}_i\rvert, \qquad
R^2 = 1-\frac{\sum_i (y_i-\hat{y}_i)^2}{\sum_i (y_i-\bar{y})^2}
\]

\paragraph{Résultats.} Le tableau~\ref{tab:resp80} et la figure~\ref{fig:r2comp} présentent les performances sur le jeu de test.

\begin{table}[H]
\centering
\caption{Performance des modèles sur le jeu de test (432 points, 26--30 juillet 2026)}
\label{tab:resp80}
\small\setlength{\tabcolsep}{5pt}
\begin{tabular}{llcccc}
\toprule
\textbf{Scénario} & \textbf{Modèle} & $\boldsymbol{R^2}$ & \textbf{RMSE (µm)} & \textbf{MAE (µm)} & \textbf{MAPE} \\
\midrule
\multirow{4}{*}{A -- en ligne}
 & Moyenne (référence) & $-0{,}025$ & 3,916 & 3,542 & 2,50\,\% \\
 & Régression Ridge & 0,890 & 1,281 & 1,002 & 0,70\,\% \\
 & Random Forest & 0,953 & 0,842 & 0,672 & 0,47\,\% \\
 & \textbf{XGBoost} & \textbf{0,950} & \textbf{0,869} & \textbf{0,669} & \textbf{0,47\,\%} \\
\midrule
\multirow{4}{*}{B -- + alimentation}
 & Moyenne (référence) & $-0{,}025$ & 3,916 & 3,542 & 2,50\,\% \\
 & Régression Ridge & 0,950 & 0,861 & 0,662 & 0,47\,\% \\
 & Random Forest & 0,952 & 0,845 & 0,667 & 0,47\,\% \\
 & \textbf{XGBoost} & \textbf{0,953} & \textbf{0,841} & \textbf{0,657} & \textbf{0,46\,\%} \\
\bottomrule
\end{tabular}
\end{table}

\begin{figure}[H]
\centering
\includegraphics[width=0.85\textwidth]{models_r2_comparison}
\caption{Coefficient de détermination $R^2$ des quatre modèles sur le jeu de test, pour les deux scénarios}
\label{fig:r2comp}
\end{figure}

La figure~\ref{fig:p80ts} superpose le P80 mesuré et le P80 estimé par XGBoost sur toute la période de test, et la figure~\ref{fig:parity} compare les valeurs point par point.

\begin{figure}[H]
\centering
\includegraphics[width=\textwidth]{p80_timeseries_B}
\caption{P80 mesuré et P80 estimé par XGBoost sur le jeu de test (scénario B)}
\label{fig:p80ts}
\end{figure}

\begin{figure}[H]
\centering
\includegraphics[width=0.45\textwidth]{p80_parity_A}\hfill
\includegraphics[width=0.45\textwidth]{p80_parity_B}
\caption{P80 estimé en fonction du P80 mesuré (la diagonale correspond à une estimation parfaite)}
\label{fig:parity}
\end{figure}

\paragraph{Robustesse.} Un seul découpage peut être favorable par hasard. Une validation croisée temporelle en 5 plis (\textit{TimeSeriesSplit}), où chaque pli teste le modèle sur une période postérieure à celle d'entraînement, donne pour XGBoost :
\begin{itemize}
  \item scénario A : $R^2 = 0{,}938 \pm 0{,}008$ et RMSE $= 0{,}93 \pm 0{,}06$~µm ;
  \item scénario B : $R^2 = 0{,}944 \pm 0{,}007$ et RMSE $= 0{,}88 \pm 0{,}03$~µm.
\end{itemize}
La faible dispersion d'un pli à l'autre montre que le modèle reste stable dans le temps.

\paragraph{Variables les plus influentes.} La figure~\ref{fig:importance} montre que, dans le scénario B, le modèle s'appuie d'abord sur la caractérisation de l'alimentation : P80 d'alimentation (55\,\% du gain total), teneur BPL (21\,\%) et débit solide d'alimentation (13\,\%). Dans le scénario A, en l'absence de ces analyses, il reconstitue l'information à partir des mesures hydrauliques : fraction solide de la surverse (36\,\%), débit de décharge du broyeur (23\,\%), débit de surverse (16\,\%) et débit d'alimentation (15\,\%). Ce comportement est cohérent avec la physique de la classification : la densité et le débit de la surverse dépendent directement de la coupure granulométrique des hydrocyclones~\cite{wills2016}.

\begin{figure}[H]
\centering
\includegraphics[width=0.75\textwidth]{xgb_feature_importance}
\caption{Importance des variables du modèle XGBoost (scénario B, gain normalisé)}
\label{fig:importance}
\end{figure}

\paragraph{Analyse.} Quatre enseignements se dégagent.
\begin{enumerate}
  \item \textbf{Le capteur virtuel est fiable.} Avec une erreur moyenne de 0,66~µm, soit 0,46\,\% du P80, XGBoost réduit l'erreur de 81\,\% par rapport à la référence constante (MAE de 3,54~µm).
  \item \textbf{Les capteurs en ligne suffisent.} Sans aucune analyse de laboratoire (scénario A), XGBoost atteint déjà $R^2 = 0{,}950$. Un modèle linéaire plafonne à 0,890 dans ce scénario : la relation entre les débits et la granulométrie n'est pas linéaire, ce qui justifie un modèle à base d'arbres. À l'inverse, dans le scénario B, la régression Ridge atteint 0,950, presque autant que XGBoost (0,953) : lorsque la caractérisation de l'alimentation est connue, un modèle linéaire suffit. L'intérêt de XGBoost se situe donc dans le cas le plus réaliste, celui où l'on ne dispose que des capteurs en ligne.
  \item \textbf{XGBoost et Random Forest sont à égalité en précision.} Les écarts, de l'ordre de 0,003 en $R^2$, ne sont pas significatifs. XGBoost se distingue par sa \textbf{sobriété} : son modèle final comporte 49 arbres de profondeur 3 et pèse \textbf{0,06~Mo}, contre 100 arbres de profondeur 22 et \textbf{8,3~Mo} pour le Random Forest, soit un modèle environ 130 fois plus léger. C'est un atout pour un déploiement embarqué, au plus près des machines.
  \item \textbf{Le $R^2$ reflète surtout la reconnaissance du type de minerai.} Les figures~\ref{fig:p80ts} et~\ref{fig:parity} montrent que le P80 évolue selon trois régimes (environ 138, 141,5 et 147~µm), qui changent en moyenne toutes les 3,3~heures et suivent la nature du minerai alimenté (P80 et teneur BPL de l'alimentation). Le modèle identifie chaque changement de régime sans retard. En revanche, \textbf{à l'intérieur d'un régime}, il n'explique pas les fluctuations : le $R^2$ calculé sur les écarts à la moyenne de chaque régime est proche de zéro ($-0{,}16$). Ces fluctuations ne sont donc liées à aucune variable d'entrée et se comportent comme un bruit de mesure, ce qui explique aussi que les trois modèles non linéaires convergent vers la même RMSE, environ 0,84~µm. Il serait donc trompeur de présenter le $R^2$ de 0,95 comme une capacité à suivre chaque micro-variation : sa valeur opérationnelle tient à la détection immédiate des changements de minerai, évaluée ci-dessous.
\end{enumerate}

\paragraph{Reproductibilité.} L'ensemble de l'expérience tient dans un script unique du projet, \nolinkurl{ml/soft_sensor_p80/train_xgboost.py}, avec une graine aléatoire fixée. Le script produit le fichier \texttt{metrics.json}, les figures de cette section et le modèle entraîné (\nolinkurl{xgb_soft_sensor_p80.joblib}).

\paragraph{Comparaison avec la pratique actuelle : l'analyse de laboratoire.} Un $R^2$ ne dit pas si l'outil est utile à l'exploitant. La vraie question est : \textit{le capteur virtuel fait-il mieux que ce dont l'opérateur dispose aujourd'hui ?} Sans granulomètre en ligne, le P80 est connu par une analyse de laboratoire périodique. Nous avons simulé cette pratique sur la même période de test : un prélèvement toutes les 4 ou 8 heures, un résultat disponible une heure après le prélèvement, et, entre deux résultats, l'opérateur se fie à la dernière valeur connue. Le capteur virtuel est ici le modèle XGBoost du scénario A, qui n'utilise que les capteurs en ligne (tableau~\ref{tab:labo}, figures~\ref{fig:labo} et~\ref{fig:labomae}).

\begin{table}[H]
\centering
\caption{Capteur virtuel et analyse de laboratoire sur le jeu de test (432 points, 26--30 juillet 2026)}
\label{tab:labo}
\small
\begin{tabularx}{\textwidth}{>{\raggedright\arraybackslash}X ccc}
\toprule
 & \textbf{Labo / 8 h} & \textbf{Labo / 4 h} & \textbf{Capteur virtuel} \\
\midrule
Erreur absolue moyenne & 4,26~µm & 3,26~µm & \textbf{0,66~µm} \\
Erreur maximale & 11,2~µm & 10,9~µm & \textbf{4,4~µm} \\
Part du temps avec une erreur $> 2$~µm & 61\,\% & 50\,\% & \textbf{2\,\%} \\
Régime de minerai correctement identifié & 39\,\% & 52\,\% & \textbf{98,6\,\%} \\
\bottomrule
\end{tabularx}
\end{table}

\begin{figure}[H]
\centering
\includegraphics[width=\textwidth]{p80_capteur_vs_labo}
\caption{48 premières heures du test : P80 réel, capteur virtuel avec son intervalle à 90\,\%, et valeur connue de l'opérateur avec une analyse de laboratoire toutes les 8~h}
\label{fig:labo}
\end{figure}

\begin{figure}[H]
\centering
\includegraphics[width=0.62\textwidth]{p80_mae_labo_vs_capteur}
\caption{Erreur absolue moyenne sur le P80 selon la source d'information}
\label{fig:labomae}
\end{figure}

Le résultat est net : l'erreur est \textbf{divisée par 6,5} par rapport à une analyse toutes les 8~heures, et par 5 par rapport à une analyse toutes les 4~heures. La raison est visible sur la figure~\ref{fig:labo} : le minerai change de nature toutes les 3~heures environ, plus vite que le laboratoire ne peut le suivre. L'opérateur pilote alors une bonne partie du temps avec une valeur périmée, souvent celle d'un autre régime. Le capteur virtuel, lui, se recale à chaque mesure, soit toutes les 15~minutes. Sur les 432 points du test, il ne se trompe de régime que 6 fois, toujours au moment même d'une transition.

\paragraph{Intervalle de confiance de chaque estimation.} Un opérateur doit savoir quand faire confiance à une estimation. Nous avons donc associé à chaque valeur un \textbf{intervalle de prédiction à 90\,\%}, construit par prédiction conforme (\textit{split conformal prediction})~\cite{angelopoulos2023}. Le principe est simple : on mesure les erreurs du modèle sur le jeu de validation, qu'il n'a pas vu pendant l'apprentissage, et l'on retient l'erreur que 90\,\% d'entre elles ne dépassent pas. On obtient un intervalle de \textbf{$\pm$\,1,41~µm}. Sur le jeu de test, \textbf{91,4\,\%} des valeurs réelles tombent dans cet intervalle, conformément à la garantie annoncée de 90\,\%. Le capteur virtuel peut ainsi afficher, par exemple, «~P80 = 147,2~µm $\pm$ 1,4~µm~».

\paragraph{Reproductibilité.} Ces analyses sont produites par le script \nolinkurl{ml/soft_sensor_p80/analyse_avancee.py}.



\subsection{Simulateur What-If}
\paragraph{Un protocole d'évaluation biaisé.} Le backend évalue le simulateur sur 20\,\% des lignes tirées au hasard dans l'historique fusionné à la minute. Or, dans cette table, chaque mesure procédé (relevée toutes les 15 minutes) est recopiée sur 15 lignes consécutives (section~\ref{sec:asymetrie}). Un tirage aléatoire place donc dans le test des lignes presque identiques à des lignes d'entraînement : le modèle est interrogé sur ce qu'il a déjà vu. Pour mesurer la précision réelle, nous avons comparé trois protocoles, avec exactement les mêmes hyperparamètres que le backend (tableau~\ref{tab:whatif}, figure~\ref{fig:whatifprot}) :
\begin{enumerate}
  \item aléatoire, à la minute : le protocole initial ;
  \item aléatoire, au pas de 15 minutes : une seule ligne par mesure procédé, ce qui supprime les doublons ;
  \item \textbf{chronologique}, au pas de 15 minutes : le modèle apprend sur les 24 premiers jours et il est testé sur les 6 derniers (25--30 juillet). C'est le seul protocole qui reproduit l'usage réel : prévoir une période que le modèle n'a jamais vue.
\end{enumerate}

\begin{table}[H]
\centering
\caption{Précision du simulateur What-If selon le protocole d'évaluation (23 sorties)}
\label{tab:whatif}
\small
\begin{tabularx}{\textwidth}{>{\raggedright\arraybackslash}X ccc}
\toprule
 & \textbf{Aléatoire, 1 min} & \textbf{Aléatoire, 15 min} & \textbf{Chronologique, 15 min} \\
\midrule
Débits, teneurs, P80 cyclones et indicateurs de circuit (14 sorties) & 0,90 -- 1,00 & 0,39 -- 1,00 & 0,33 -- 1,00 \\
\quad dont P80 de la surverse & 0,991 & 0,937 & \textbf{0,956} \\
\quad dont P80 de décharge du broyeur & 0,903 & 0,385 & \textbf{0,329} \\
Signaux de santé et pressions (9 sorties) & 0,58 -- 0,78 & $-0{,}10$ -- 0,02 & $\mathbf{-0{,}06}$ \textbf{--} $\mathbf{0{,}02}$ \\
\midrule
$R^2$ médian (23 sorties) & 0,996 & 0,976 & \textbf{0,970} \\
Sorties avec $R^2 > 0{,}9$ & 14 & 13 & \textbf{13} \\
Sorties avec $R^2 < 0{,}5$ & 0 & 10 & \textbf{10} \\
\bottomrule
\end{tabularx}
\end{table}

\begin{figure}[H]
\centering
\includegraphics[width=0.92\textwidth]{whatif_r2_protocoles}
\caption{$R^2$ de chaque sortie du simulateur What-If : protocole initial (gris) et protocole chronologique (bleu)}
\label{fig:whatifprot}
\end{figure}

\paragraph{Analyse.} Deux conclusions nettes se dégagent.
\begin{enumerate}
  \item \textbf{Le simulateur est fiable pour les grandeurs de procédé.} Les débits, les teneurs BPL, les P80 des cyclones, la charge circulante et le P80 de la surverse conservent un $R^2$ compris entre 0,956 et 0,999 sur une période future. Ces sorties sont gouvernées par les bilans matière du circuit, que les 8 entrées décrivent directement. Seul le P80 de décharge du broyeur reste mal prédit (0,33), faute d'information sur l'état de la charge broyante.
  \item \textbf{Les signaux de santé ne sont pas prévisibles à partir des réglages.} Les scores de 0,58 à 0,78 du protocole initial disparaissent dès que les doublons sont supprimés : ils provenaient uniquement de la fuite d'information. C'est physiquement logique : la puissance, les températures de paliers, les vibrations et les pressions évoluent avec l'\textbf{usure} des équipements et le temps écoulé depuis la dernière intervention, pas avec la consigne de débit ou de vitesse choisie par l'opérateur.
\end{enumerate}
Cette évaluation conduit à une recommandation claire : le simulateur What-If doit être réservé aux 13 sorties de procédé fiables, et la prévision des signaux de santé doit être confiée à un modèle de dégradation, présenté dans la section suivante.

\subsection{Maintenance prédictive : pronostic de durée de vie résiduelle}
\label{sec:pronostic}

\paragraph{Objectif.} Le tableau de bord de maintenance affiche l'état actuel des équipements. Un jumeau numérique doit aller plus loin et répondre à la question de l'exploitant : \textit{dans combien de jours cet équipement devra-t-il être arrêté ?} C'est l'estimation de la \textbf{durée de vie résiduelle} (\textit{Remaining Useful Life}, RUL), cœur de la maintenance conditionnelle et prévisionnelle~\cite{jardine2006}.

\paragraph{Démarche.} Nous suivons un indicateur d'usure par équipement critique : la vibration de la pompe SP\_001, la vibration du broyeur BM\_001 et l'indice d'usure des apex des hydrocyclones CY\_001. Les signaux de santé, enregistrés chaque minute, sont moyennés par heure. L'analyse se déroule en quatre étapes :
\begin{enumerate}
  \item détecter automatiquement les interventions dans les signaux et les rapprocher de la GMAO ;
  \item modéliser la dégradation entre deux interventions ;
  \item valider le pronostic a posteriori, sur les cycles déjà terminés ;
  \item pronostiquer la date à laquelle chaque indicateur atteindra son seuil d'alarme.
\end{enumerate}

\paragraph{Étape 1 -- Détection des interventions et audit de la GMAO.} Une intervention de maintenance se traduit par une chute brutale de l'indicateur d'usure : la machine repart «~comme neuve~». Nous considérons comme intervention toute baisse horaire supérieure à 8 fois l'écart-type du bruit, estimé de façon robuste par l'écart absolu médian. Cinq interventions sont détectées sur le mois, puis rapprochées des ordres de travail de la GMAO (tableau~\ref{tab:interventions}).

\begin{table}[H]
\centering
\caption{Interventions détectées dans les signaux et rapprochement avec la GMAO}
\label{tab:interventions}
\small
\begin{tabularx}{\textwidth}{l l l X}
\toprule
\textbf{Équipement} & \textbf{Indicateur} & \textbf{Date détectée} & \textbf{Ordre de travail dans la GMAO} \\
\midrule
SP\_001 (pompe) & Vibration & 10/07/2026 & \textbf{Aucun} \\
SP\_001 (pompe) & Vibration & 23/07/2026 & WO-2026-0814 : remplacement du \textit{throatbush} et du revêtement de la roue (correctif) \\
BM\_001 (broyeur) & Vibration & 15/07/2026 & \textbf{Aucun} \\
BM\_001 (broyeur) & Vibration & 28/07/2026 & WO-2026-0820 : alerte de température du palier, filtre et lubrifiant (conditionnel) \\
CY\_001 (cyclones) & Usure des apex & 19/07/2026 & Trois ordres de travail : remplacement des apex et des \textit{vortex finders} des batteries A et B, débouchage de la batterie C \\
\bottomrule
\end{tabularx}
\end{table}

Trois interventions sur cinq correspondent exactement, au jour près, à un ordre de travail. Les deux autres (10/07 sur la pompe, 15/07 sur le broyeur) sont clairement visibles dans les signaux mais \textbf{absentes de la GMAO}. Le jumeau numérique sert ici d'outil d'audit : il révèle des interventions non tracées, qui faussent ensuite les indicateurs de fiabilité calculés à partir de la GMAO (MTBF, MTTR). Une autre incohérence apparaît : l'ordre WO-2026-0820 mentionne une alerte à 68~°C sur un palier du broyeur, alors que les deux capteurs de température de palier du jumeau n'ont jamais dépassé 62,2~°C. L'alerte provient donc probablement d'un capteur qui n'est pas encore remonté dans le jumeau.

\paragraph{Étape 2 -- Modèle de dégradation.} Entre deux interventions, chaque indicateur croît de façon quasi linéaire (tableau~\ref{tab:cycles}). Surtout, la \textbf{vitesse de dégradation est la même d'un cycle à l'autre} : 0,045~mm/s par jour pour la pompe, 0,060~mm/s par jour pour le broyeur et 2,60~\% par jour pour les apex. C'est ce qui rend le pronostic possible : connaître la pente et le niveau actuel suffit pour prévoir la suite.

\begin{table}[H]
\centering
\caption{Cycles de dégradation observés (entre deux interventions)}
\label{tab:cycles}
\small\setlength{\tabcolsep}{4pt}
\begin{tabular}{llcccc}
\toprule
\textbf{Équipement} & \textbf{Cycle} & \textbf{Durée} & \textbf{Début $\rightarrow$ fin} & \textbf{Pente / jour} & \textbf{$R^2$} \\
\midrule
Pompe (mm/s) & 01/07 -- 09/07 & 9,0 j & 1,80 $\rightarrow$ 2,20 & 0,0448 & 0,992 \\
 & 10/07 -- 22/07 & 13,0 j & 1,81 $\rightarrow$ 2,38 & 0,0452 & 0,996 \\
 & 23/07 -- en cours & 8,0 j & 1,80 $\rightarrow$ 2,15 & 0,0453 & 0,990 \\
\midrule
Broyeur (mm/s) & 01/07 -- 14/07 & 14,0 j & 2,61 $\rightarrow$ 3,43 & 0,0595 & 0,996 \\
 & 15/07 -- 27/07 & 13,0 j & 2,60 $\rightarrow$ 3,37 & 0,0600 & 0,996 \\
 & 28/07 -- en cours & 3,0 j & 2,60 $\rightarrow$ 2,77 & 0,0604 & 0,932 \\
\midrule
Apex (\%) & 01/07 -- 18/07 & 18,0 j & 5,3 $\rightarrow$ 51,5 & 2,600 & 1,000 \\
 & 19/07 -- en cours & 12,0 j & 5,3 $\rightarrow$ 35,9 & 2,600 & 1,000 \\
\bottomrule
\end{tabular}
\end{table}

\paragraph{Étape 3 -- Validation a posteriori (\textit{backtest}).} Pour vérifier le pronostic sans attendre l'avenir, nous nous replaçons dans le passé. À chaque heure d'un cycle terminé, le modèle est ajusté \textbf{uniquement sur les données disponibles à cet instant}, puis il prédit la date à laquelle l'indicateur atteindra le niveau auquel l'intervention a réellement eu lieu. On compare ensuite cette date à la date réelle. Les cinq cycles terminés fournissent ainsi 1\,543 pronostics (tableau~\ref{tab:backtest}, figure~\ref{fig:backtest}).

\begin{table}[H]
\centering
\caption{Erreur du pronostic de durée de vie résiduelle selon l'horizon (5 cycles, 1\,543 pronostics)}
\label{tab:backtest}
\small
\begin{tabular}{lcc}
\toprule
\textbf{Moment du pronostic} & \textbf{Erreur moyenne} & \textbf{Erreur maximale} \\
\midrule
7 jours avant l'intervention & 4,6~h & 25,2~h \\
3 jours avant l'intervention & 2,4~h & 4,8~h \\
1 jour avant l'intervention & 2,5~h & 4,7~h \\
\midrule
Ensemble des pronostics & 6,4~h (6,0\,\% de la durée restante) & -- \\
\bottomrule
\end{tabular}
\end{table}

\begin{figure}[H]
\centering
\includegraphics[width=0.6\textwidth]{pronostic_backtest}
\caption{Durée de vie résiduelle prédite en fonction de la durée réelle ; la diagonale correspond à un pronostic parfait}
\label{fig:backtest}
\end{figure}

Une semaine avant l'intervention, la date est prédite à moins de 5~heures près en moyenne. Trois jours avant, l'erreur ne dépasse jamais 5~heures : c'est largement suffisant pour planifier une équipe, commander une pièce ou caler l'arrêt sur un changement de poste. La figure~\ref{fig:backtest} montre aussi la limite de la méthode : pendant les 1 à 2 premiers jours d'un cycle, la pente est estimée sur trop peu de points et le pronostic oscille. Il ne doit donc être affiché qu'après environ 48~heures de fonctionnement depuis la dernière intervention. Enfin, 90\,\% des pronostics ont une erreur inférieure à 8\,\% de la durée restante : nous utilisons cette valeur pour construire l'intervalle de confiance des pronostics de l'étape suivante.

\paragraph{Étape 4 -- Pronostic à la fin des données.} Le modèle est enfin appliqué au cycle en cours de chaque équipement, au 30 juillet 2026. Les seuils d'alarme proviennent de la norme ISO 10816-3 pour les vibrations~\cite{iso10816} (limite entre la zone B, «~fonctionnement acceptable sans restriction~», et la zone C, «~fonctionnement acceptable seulement pour une durée limitée~»), et de la pratique observée dans la GMAO pour les apex (tableau~\ref{tab:pronostic}, figure~\ref{fig:pronostic}).

\begin{table}[H]
\centering
\caption{Pronostic au 30 juillet 2026 (intervalle à 90\,\%)}
\label{tab:pronostic}
\small
\begin{tabularx}{\textwidth}{l c >{\raggedright\arraybackslash}X c c}
\toprule
\textbf{Équipement} & \textbf{Niveau actuel} & \textbf{Seuil d'alarme} & \textbf{Durée restante} & \textbf{Date prévue} \\
\midrule
Pompe SP\_001 & 2,15~mm/s & 2,8~mm/s (ISO 10816-3, groupe 2, zones B/C) & 14,1 j [13,0 -- 15,2] & 14/08/2026 \\
Broyeur BM\_001 & 2,77~mm/s & 4,5~mm/s (ISO 10816-3, groupe 1, zones B/C) & 28,5 j [26,2 -- 30,8] & 28/08/2026 \\
Apex CY\_001 & 35,9~\% & 50~\% (niveau de remplacement du 19/07) & 5,3 j [4,9 -- 5,7] & 05/08/2026 \\
\bottomrule
\end{tabularx}
\end{table}

\begin{figure}[H]
\centering
\includegraphics[width=\textwidth]{pronostic_equipements}
\caption{Indicateurs d'usure, interventions détectées (trait plein : présente dans la GMAO ; pointillé : absente) et pronostic jusqu'au seuil d'alarme (zone grisée)}
\label{fig:pronostic}
\end{figure}

\paragraph{Analyse.} Ce pronostic apporte trois enseignements concrets.
\begin{enumerate}
  \item \textbf{La prochaine intervention urgente concerne les hydrocyclones.} Les apex atteindront le niveau de remplacement vers le 5 août, soit dans 5 jours. C'est l'information la plus utile pour l'exploitant : elle laisse le temps de préparer les pièces et de planifier l'arrêt.
  \item \textbf{Les interventions pourraient être espacées.} La pompe a été arrêtée en moyenne tous les 11 jours, à une vibration de 2,29~mm/s, alors qu'au rythme de dégradation observé elle n'atteindrait la zone C qu'après 22 jours environ. Pour le broyeur, les interventions ont eu lieu tous les 13,5 jours, contre environ 32 jours avant d'atteindre son seuil. Selon ce seul critère vibratoire, les équipements sont donc arrêtés à mi-vie environ. Avant de modifier la stratégie, il faudra toutefois confirmer cette marge avec les autres critères de décision des techniciens, en particulier l'épaisseur résiduelle des revêtements d'usure.
  \item \textbf{La GMAO doit être fiabilisée.} Deux interventions sur cinq n'y figurent pas. Le jumeau permet de les détecter automatiquement et pourrait proposer de créer l'ordre de travail manquant.
\end{enumerate}

\paragraph{Limites.} Dans les données simulées, la dégradation est linéaire et très régulière, ce qui explique la précision obtenue. Sur une installation réelle, l'usure s'accélère souvent en fin de vie, et un modèle exponentiel ou par morceaux serait alors nécessaire. Le mois d'historique ne contient en outre que cinq cycles complets. Ces résultats valident donc la démarche et l'outillage, mais leur précision devra être confirmée sur des mesures réelles. L'ensemble est produit par le script \nolinkurl{ml/maintenance_predictive/pronostic.py}.


\subsection{Persistance des données}
L'historien a été vérifié par des tests automatisés qui simulent un arrêt complet (tableau~\ref{tab:histtests}). La suite de tests du backend (15 tests) passe entièrement.

\begin{table}[H]
\centering
\caption{Vérification de l'historien de télémétrie}
\label{tab:histtests}
\begin{tabularx}{\textwidth}{X c}
\toprule
\textbf{Scénario de test} & \textbf{Résultat} \\
\midrule
251 messages reçus, arrêt, puis redémarrage de l'historien : les 251 lignes sont retrouvées & Réussi \\
Après redémarrage, les 100 derniers messages de chaque tag sont restaurés en mémoire & Réussi \\
Démarrage de l'API, réception de 30 messages Node-RED, arrêt complet puis redémarrage : 30 lignes conservées, historique restauré & Réussi \\
Requête filtrée par tag et export CSV par l'API & Réussi \\
Valeur non numérique : conservée dans le message JSON, sans erreur & Réussi \\
\bottomrule
\end{tabularx}
\end{table}

\subsection{Vérification des besoins}
\begin{table}[H]
\centering
\caption{Couverture des besoins fonctionnels}
\label{tab:couverture}
\begin{tabular}{llc}
\toprule
\textbf{Réf.} & \textbf{Besoin} & \textbf{Statut} \\
\midrule
BF1 & Authentification et rôles & Réalisé \\
BF2 & Acquisition (relecture + OPC UA/Node-RED) & Réalisé \\
BF3 & Pilotage de la simulation & Réalisé \\
BF4 & Synoptique et fiches équipements & Réalisé \\
BF5 & Graphe de connaissances & Réalisé \\
BF6 & Vue 3D & Réalisé \\
BF7 & Maintenance et fiabilité & Réalisé \\
BF8 & Simulateur What-If & Réalisé \\
BF9 & Copilote IA & Réalisé \\
BF10 & Historique persistant & Réalisé \\
\bottomrule
\end{tabular}
\end{table}

\section{Limites}
\label{sec:limites}
Une évaluation honnête du prototype conduit à relever plusieurs limites :
\begin{itemize}
  \item \textbf{Données simulées.} Les deux sources proviennent de données simulées (qualité \texttt{SIM} dans les messages). Les performances des modèles devront être confirmées sur des mesures réelles, plus bruitées et plus lacunaires.
  \item \textbf{Historique court.} Un mois de données seulement : la période de test du capteur virtuel couvre 4,5 jours, et le pronostic repose sur cinq cycles de dégradation complets. Les dérives lentes (saisons, changement de gisement) ne sont pas représentées.
  \item \textbf{What-If limité aux grandeurs de procédé.} L'évaluation chronologique montre que le simulateur ne sait pas prédire les signaux de santé à partir des réglages. L'interface affiche encore ces sorties, avec le score issu du protocole aléatoire initial : elle devra être alignée sur les résultats de ce chapitre.
  \item \textbf{Absence de boucle de retour automatique.} Au sens de Kritzinger \textit{et al.}~\cite{kritzinger2018}, le système est une ombre numérique : il n'agit pas sur le procédé.
  \item \textbf{Pas de modèle physique.} Les modèles sont orientés données. Ils n'extrapolent pas en dehors des régimes observés dans l'historique.
  \item \textbf{Dépendance à un service externe.} Le copilote privilégie une API de LLM hébergée, ce qui pose des questions de confidentialité en milieu industriel. Le moteur local de secours n'en atténue qu'en partie l'impact.
\end{itemize}

\section{Conclusion}
Le prototype couvre l'ensemble des besoins fonctionnels identifiés, de l'acquisition OPC UA jusqu'à l'aide à la décision. Les trois briques d'IA ont été évaluées avec des protocoles qui reproduisent leur usage réel. Le capteur virtuel divise par 6,5 l'erreur commise avec une analyse de laboratoire toutes les 8~heures, et il accompagne chaque estimation d'un intervalle de confiance vérifié. Le simulateur What-If est fiable sur les 13 grandeurs de procédé qu'il doit prévoir. Le module de pronostic prédit la date d'intervention à quelques heures près, trois jours à l'avance, et il a révélé deux interventions absentes de la GMAO. L'historique survit désormais aux redémarrages, et les limites relevées tracent directement la feuille de route présentée en conclusion.
